\documentclass[a4paper,10pt]{report}

\def\packagepath{../../preambule}   % path du package princial
\usepackage{\packagepath/preambule}    % utilisation du fichier de configuration

\def\level{Première }              % Classe
\def\course{NSI}              % Matière
\def\eval{DS05}

\def\visibleornot{visible}    % visible or invisible
\def\documentpath{./Sources_Latex}
\def\date{Jeudi 22 Novembre 2025}
\renewcommand{\arraystretch}{1}  % Ecart dans les tableau

\def\ptsexoA{4}
\def\ptsexoB{2}
\def\ptsexoC{2}
\def\ptsexoD{2}
\def\ptsexoE{1}
\def\ptsexoF{4}
\def\ptsexoG{2}
\def\ptsexoH{1}
\def\ptsexoI{2}
\def\ptsexoJ{2}
\def\ptsexoK{1}

\def\ptstotal{\ptsexoA+\ptsexoB}

\usepackage{hyperref}
\usepackage{tasks}

\usetikzlibrary{shapes, arrows, positioning}
\usetikzlibrary{shapes.geometric, arrows, positioning, decorations.pathreplacing}

\tikzstyle{etat} = [draw, rounded corners=15pt, minimum width=2.5cm, minimum height=1.2cm, text centered, font=\bfseries]
\tikzstyle{nouveau} = [etat, fill=blue!40]
\tikzstyle{pret} = [etat, fill=green!60]
\tikzstyle{elu} = [etat, fill=yellow!60]
\tikzstyle{bloque} = [etat, fill=red!60]
\tikzstyle{termine} = [etat, fill=white]
\tikzstyle{fleche} = [thick, ->, >=stealth]

\usepackage{float} % Permet d'utiliser [H]
\usepackage[T1]{fontenc}
\begin{document}

%%%%%%%%%%%%%%%%%%%%%%%%%%%%%%%%%%%%%%%%%%%%%%%%%%ù

%\PageGardeSujetBac[Matiere = Numérique et Sciences Informatiques,
%                  Session = 2025,
%                  AffJour=false,
%                  Duree = {3 heures 30},
%                  DernierePage = \pageref{LastPage},
%                  ModeExamen=false
%                  ]
\renewcommand{\labelitemi}{\textbullet} %pour éviter les tirets dans les "itemize" qui apportent confusion avec le signe moins.
%% Début page de garde style BAC

   %%%%%%%%%%%%%%%%%%%%%%%%%%%%%%%%%%%%%%%%%%%%%%%%%%%%%%%%
%\PageGardeSujetBac[clés]

\pagestyle{DS_FP}          % Feuille de style fancy
\NomPrenomNote{}

\emph{Dans tout le devoir, sauf indication contraire, on considèrera que les nombres sont des entiers relatifs codés sur 8 bits exactement. On supposera que les variables d'entrée de chaque algorithme sont conformes aux attentes de l'algorithme.
}
\bigskip

\exods{\ptsexoA}

\begin{enumerate}
    \item Sans aucune justification, donner le nombre d'états possibles codables avec:
    
\begin{enumerate}[a)]
    \item 3 bits \dotfill{}
    \item 8 bits \dotfill{}
    \item n bits \dotfill{}
\end{enumerate}

\item On code maintenant un nombre entier relatif avec la méthode du complément à 2 vue en classe. Quelle est la valeur minimale et maximale décimale codable avec:

\begin{enumerate}[a)]
    \item 3 bits \dotfill{}
    \item 8 bits \dotfill{}
    \item n bits \dotfill{}
\end{enumerate}

\item Que peut-on dire de l'entier relatif codé sur 8 bits en complément à deux suivant: \verb|1001 0101|. Sa valeur décimale est-elle: (plusieurs réponses possibles)

\begin{tasks}(4)
    \task[$\bigcirc$] paire
    \task[$\bigcirc$] impaire
    \task[$\bigcirc$] positive
    \task[$\bigcirc$] négative
\end{tasks}

\end{enumerate}


\medskip

\begin{minipage}{0.49\linewidth}
    \exods{\ptsexoB}

 \'Ecrire l'algorithme de conversion d'un entier naturel compris entre 0 et 255 en binaire sur 8 bits en pseudo code ou en Python.

 \begin{answer}{1}{6} 
    \begin{verbatim}
    \end{verbatim}   
\end{answer}
\end{minipage}\hfill{}
\begin{minipage}{0.49\linewidth}
    \exods{\ptsexoC}

    Citez les différentes étapes de la conversion d'un entier relatif en binaire sur 8 bits. Vous choisirez de développer un des deux algorithmes vus en cours. Inutile de l'écrire en Python.
    
    \begin{answer}{1}{6} 
    \begin{verbatim}
    \end{verbatim}   
\end{answer}
\end{minipage}



\medskip



\begin{minipage}{0.49\linewidth}
 \exods{\ptsexoD}

Convertir le nombre entier relatif suivant avec la méthode du complément à 2 sur 8 bits. Justifier.
\begin{answer}{1}{5} 
    \begin{center}
        \verb|-56|
    \end{center}   
\end{answer}
\end{minipage}\hfill
\begin{minipage}{0.49\linewidth}

     \exods{\ptsexoE}

Convertir le nombre entier relatif suivant avec la méthode du complément à 2 sur 8 bits. Justifer.
\begin{answer}{1}{5} 
    \begin{center}
        \verb|112|
    \end{center}   
\end{answer}
\end{minipage}

\medskip

\pagebreak
\thispagestyle{DS_LP}
\exods{\ptsexoF}

On suppose que l'on dispose des fonctions suivantes en Python (il est inutile d'écrire le code de ces fonctions). Ecrire la fonction \texttt{conversion\_relatif\_binaire} qui convertit un entier relatif en binaire sur 8 bits. Vous utiliserez les fonctions précédentes pour créer votre fonction. 
 Vous appliquerez la méthode de votre choix (la méthode 1 est plus simple !)


\begin{CodePythontexAlt}[TaillePolice=\large, Largeur=1\linewidth,Centre,PremLigne=1]{}
def conversion_entier_binaire(val_entiere):
    ''' Convertit un entier naturel en binaire sur 8 bits. 
    val_entiere : int : entier naturel compris entre 0 et 255
    val_binaire : str : représentation binaire sur 8 bits de l'entier naturel'''
    # Code non demandé
    return val_binaire    
assert conversion_entier_binaire(34) == '00100010'

def inverse_bit(bit):
    ''' Inverse un bit 
    bit : str : bit à inverser
    bit_inverse : str : bit inversé'''
    # Code non demandé
    return bit_inverse
assert inverse_bit('0') == '1'
assert inverse_bit('1') == '0'

def additionne_binaire(val_binaire):
    ''' Ajoute 1 à une valeur binaire 
    val_binaire : str : mot de 8 bits
    somme_binaire : str : mot de 8 bits 
    Si création d'un 9eme bit, celui-ci est perdu '''
    # Code non demandé
    return somme_binaire
assert additionne_binaire('00000001') == '00000010'
assert additionne_binaire('11111111') == '00000000'

\end{CodePythontexAlt}

  
    \begin{answer}{1}{10.}
    \begin{verbatim}
def conversion_relatif_binaire(val_entiere):
    ''' Convertit un entier relatif en binaire sur 8 bits. 
    val_entiere : int : entier relatif compris entre -128 et 127
    return : str : représentation binaire sur 8 bits de l'entier relatif'''
    \end{verbatim}
    \end{answer}


\end{document}